% !PRINTED 216
\weberpage{216}

\begin{center}
\S.\,57.

\gesperrt{The Octahedral Group.}
\end{center}

In the octahedral group there occurs a system of six conjugate
four-fold poles. We accordingly have a substitution of degree 4
whose period
\[
1, \ \Theta, \ \Theta^2, \ \Theta^3, \ \Theta^4 = 1
\]
gives a pair of these conjugate four-fold poles. By \S.\,53, 1.
we may assume in the group the substitutions:
\weq{1}{\Theta(x) = ix, \quad \psi(x) = \frac{1}{x},}
and we have
\weq{2}{\psi^2 = 1, \quad \Theta\psi = \psi\Theta^3, \quad \Theta^2\psi = \psi\Theta^2, \quad \Theta^3\psi = \psi\Theta.}
The fundamental form of degree 6 belonging to the system of the
four-fold poles must remain unchanged, up to a constant factor,
under the substitutions $\Theta$ and $\psi$, and it must moreover
contain the factor $x_1x_2$, and hence be of one of the two forms
\[
x_1x_2(x_1^4 - x_2^4), \quad x_1x_2(x_1^4 + x_2^4).
\]
It is immaterial which of the two assumptions we pursue, since
the one passes into the other by the substitution of $ix_2$ for $x_2$,
which corresponds to the passage to a transformed group. If we take
\weq{3}{f(x_1, x_2) = x_1x_2(x_1^4 - x_2^4)}
as the fundamental form of degree 6, then the six four-fold poles are
\weq{4}{0, \ \infty, \ 1, \ -1, \ i, \ -i.}
Now by \S.\,53, 2. there must be in $G$ a substitution $\chi$
such that
\[
\chi(-1) = 0, \quad \chi(+1) = \infty
\]
and if one determines the constant factor as with the tetrahedral
group [\S.\,56, (2)], it is found that
\weq{5}{\chi(x) = i\,\frac{x+1}{x-1}.}
From this we obtain
\[
\chi^2(x) = \frac{x+i}{x-i}, \quad \chi^3(x) = x,
\]
%CONT>>
